\begin{proposition}
\label{proposition-lefschetz}
In the situation above assume there exists an integer $\sigma$ such that
for all points $p \in P \setminus Q$ we have
$$
\text{depth}(\mathcal{F}_p) + \dim(\overline{\{p\}}) > \sigma
$$
Then the map
$$
H^i(P, \mathcal{F}) \longrightarrow \lim H^i(Q_n, \mathcal{F}_n)
$$
is an isomorphism for $0 \leq i < \sigma$.
\end{proposition}

\begin{proof}
We will use More on Morphisms, Lemma \ref{more-morphisms-lemma-apply-proj-spec}
and we will use the notation used and results found
More on Morphisms, Section \ref{more-morphisms-section-proj-spec}
without further mention; this proof will not make sense without
at least understanding the statement of the lemma.
Observe that in our case
$A = \bigoplus_{m \geq 0} \Gamma(P, \mathcal{L}^{\otimes m})$
is a finite type $k$-algebra
all of whose graded parts are finite dimensional $k$-vector spaces, see
Cohomology of Schemes, Lemma \ref{coherent-lemma-coherent-proper-ample}.

\medskip\noindent
We may and do think of $s$ as an element $f \in A_1 \subset A$, i.e.,
a homogeneous element of degree $1$ of $A$. Denote
$Y = V(f) \subset X$ the closed subscheme defined by $f$.
Then $U \cap Y = (\pi|_U)^{-1}(Q)$ scheme theoretically.
Recall the notation
$\mathcal{F}_U = \pi^*\mathcal{F}|_U = (\pi|_U)^*\mathcal{F}$.
This is a coherent $\mathcal{O}_U$-module.
Choose a finite $A$-module $M$ such that
$\mathcal{F}_U = \widetilde{M}|_U$
(for existence see Local Cohomology, Lemma
\ref{local-cohomology-lemma-finiteness-pushforwards-and-H1-local}).
We claim that $H^i_Z(M)$ is annihilated by
a power of $f$ for $i \leq \sigma + 1$.

\medskip\noindent
To prove the claim we will apply
Local Cohomology, Proposition \ref{local-cohomology-proposition-annihilator}.
Translating into geometry we see that it suffices to prove
for $u \in U$, $u \not \in Y$ and $z \in \overline{\{u\}} \cap Z$
that
$$
\text{depth}(\mathcal{F}_{U, u}) +
\dim(\mathcal{O}_{\overline{\{u\}}, z}) > \sigma + 1
$$
This requires only a small amount of thought.

\medskip\noindent
Observe that $Z = \Spec(A_0)$ is a finite set of closed points of $X$ because
$A_0$ is a finite dimensional $k$-algebra.
(The reader who would like $Z$ to be a singleton can replace the
finite $k$-algebra $A_0$ by $k$; it won't affect anything else in the proof.)

\medskip\noindent
The morphism $\pi : L \to P$ and its restriction $\pi|_U : U \to P$
are smooth of relative dimension $1$.
Let $u \in U$, $u \not \in Y$ and $z \in \overline{\{u\}} \cap Z$.
Let $p = \pi(u) \in P \setminus Q$ be its image.
Then either $u$ is a generic
point of the fibre of $\pi$ over $p$ or a closed point of the fibre.
If $u$ is a generic point of the fibre, then
$\text{depth}(\mathcal{F}_{U, u}) = \text{depth}(\mathcal{F}_p)$
and
$\dim(\overline{\{u\}}) = \dim(\overline{\{p\}}) + 1$.
If $u$ is a closed point of the fibre, then
$\text{depth}(\mathcal{F}_{U, u}) = \text{depth}(\mathcal{F}_p) + 1$
and
$\dim(\overline{\{u\}}) = \dim(\overline{\{p\}})$.
In both cases we have
$\dim(\overline{\{u\}}) = \dim(\mathcal{O}_{\overline{\{u\}}, z})$
because every point of $Z$ is closed. Thus the desired
inequality follows from the assumption in the statement of
the lemma.

\medskip\noindent
Let $A'$ be the $f$-adic completion of $A$. So $A \to A'$ is flat by
Algebra, Lemma \ref{algebra-lemma-completion-flat}.
Denote $U' \subset X' = \Spec(A')$ the inverse image of
$U$ and similarly for $Y'$ and $Z'$. Let $\mathcal{F}'$
on $U'$ be the pullback of $\mathcal{F}_U$ and let
$M' = M \otimes_A A'$.
By flat base change for local cohomology
(Local Cohomology, Lemma \ref{local-cohomology-lemma-torsion-change-rings})
we have
$$
H^i_{Z'}(M') = H^i_Z(M) \otimes_A A'
$$
and we find that for $i \leq \sigma + 1$
these are annihilated by a power of $f$. Consider the diagram
$$
\xymatrix{
& H^i(U, \mathcal{F}_U) \ar[ld] \ar[d] \ar[r] &
\lim H^i(U, \mathcal{F}_U/f^n\mathcal{F}_U) \ar@{=}[d] \\
H^i(U, \mathcal{F}_U) \otimes_A A' \ar@{=}[r] &
H^i(U', \mathcal{F}') \ar[r] &
\lim H^i(U', \mathcal{F}'/f^n\mathcal{F}')
}
$$
The lower horizontal arrow is an isomorphism for $i < \sigma$ by
Lemma \ref{lemma-alternative-higher} and the torsion
property we just proved. The horizontal equal sign is flat base change
(Cohomology of Schemes, Lemma \ref{coherent-lemma-flat-base-change-cohomology})
and the vertical equal sign is because $U \cap Y$ and $U' \cap Y'$
as well as their $n$th infinitesimal neighbourhoods
are mapped isomorphically onto each other (as we are
completing with respect to $f$).

\medskip\noindent
Applying More on Morphisms, Equation
(\ref{more-morphisms-equation-cohomology-torsor})
we have compatible direct sum decompositions
$$
\lim H^i(U, \mathcal{F}_U/f^n\mathcal{F}_U) =
\lim
\left(
\bigoplus\nolimits_{m \in \mathbf{Z}}
H^i(Q_n, \mathcal{F}_n \otimes \mathcal{L}^{\otimes m})
\right)
$$
and
$$
H^i(U, \mathcal{F}_U) =
\bigoplus\nolimits_{m \in \mathbf{Z}}
H^i(P, \mathcal{F} \otimes \mathcal{L}^{\otimes m})
$$
Thus we conclude by Algebra, Lemma \ref{algebra-lemma-daniel-litt}.
\end{proof}